% !PRINTED 533
\weberpage{533}

First of all it is clear that a basis of a functional $\mu$ is at the
same time a basis of all the functionals associated with $\mu$; further,
that every element $\alpha_1, \alpha_2, \ldots$ of such a basis must be
divisible by $\mu$; and finally, that from one basis of $\mu$ all the
others can be derived by applying a linear substitution with integral
rational coefficients and determinant $\pm 1$.

In order now to form a basis of $\mu$, we start from a minimal basis
(basis of $\mathfrak{o}$), which, as above, we denote by
\[
\omega_1, \omega_2, \ldots, \omega_n
\]

Since there are positive integral rational numbers that are divisible by
$\mu$, there are also integral rational numbers $a$ for which the
product $a\,\omega_1$ is divisible by $\mu$. The \gesperrt{least} positive among
these numbers we will denote by $a_{1,1}$, and we set
\[
\alpha_1 = a_{1,1}\,\omega_1
\]
Thereupon we denote by $a_{2,2}$ the \gesperrt{least positive} integral rational
number for which an integral rational $a_{1,2}$ can be determined in
such a way that
\[
\alpha_2 = a_{1,2}\,\omega_1 + a_{2,2}\,\omega_2
\]
becomes divisible by $\mu$, and so we continue. In this way we form the
system
\weq{3}{
\begin{aligned}
\alpha_1 &= a_{1,1}\,\omega_1 \\
\alpha_2 &= a_{1,2}\,\omega_1 + a_{2,2}\,\omega_2 \\
&\cdots\cdots\cdots\cdots\cdots\cdots\cdots\cdots\cdots\cdots \\
\alpha_n &= a_{1,n}\,\omega_1 + a_{2,n}\,\omega_2 + \cdots + a_{n,n}\,\omega_n,
\end{aligned}
}
\gesperrt{in which $a_{1,1}, a_{2,2}, \ldots, a_{n,n}$ are the least
positive integral rational numbers that make possible a determination of
the integral rational numbers $a_{1,2}, \ldots, a_{n-1,n}$ such that the
integers $\alpha_1, \alpha_2, \ldots, \alpha_n$ become divisible by $\mu$.}

Every number $a_{r,r}$ is, according to this, defined
independently of the others by the fact that it is the least natural
number for which the associated integral rational numbers
$a_{1,r}, a_{2,r}, \ldots, a_{r-1,r}$ can be so determined that
\weq{4}{
\alpha_r = a_{1,r}\,\omega_1 + a_{2,r}\,\omega_2 + \cdots +
a_{r-1,r}\,\omega_{r-1} + a_{r,r}\,\omega_r
}
becomes divisible by $\mu$. In the particular case where
$\omega_r$ is itself divisible by $\mu$, one has accordingly to set
$a_{r,r} = 1$, and
%CONT>>
